\chapter{Improvements and Extensions of the Synthetic Generation Pipeline}
\label{chap:evaluation}

Having demonstrated the viability of our privacy-safe synthetic document generation approach on MIMIC-III through ICD-9 classification and NER tasks, we identified several opportunities to extend and refine our methodology. The previous chapter established that our keyword-guided generation coupled with reward-based refinement successfully produces high-quality synthetic documents that can match or exceed the performance of models trained on real data. However, the evaluation remained constrained to specific task types and datasets, and several limitations emerged regarding task complexity, evaluation robustness, and privacy guarantees.

This chapter explores three key extensions to address these limitations and broaden the applicability of our approach. First, we extend our evaluation to more complex conversational medical tasks using the AlpaCare benchmark, enabling direct comparison with related privacy-preserving methods like KnowledgeSG. Second, we investigate advanced filtering and evaluation techniques—including medical domain filtering and LLM-as-judge mechanisms—to improve the quality of synthetic document selection beyond the SemScore metric employed in our initial pipeline. Third, we conduct rigorous privacy analysis through linkage and discrimination attacks to quantify re-identification risks and validate the privacy guarantees of our generated documents.

To support these investigations, we build upon our core synthetic generation pipeline, which addresses the critical need for producing high-quality training data while preserving patient information confidentiality. The approach relies on decoder-type generative models (we use Llama2-7B \cite{touvron2023llama} and Mistral-7B \cite{jiang2023mistral} in our experiments) coupled with an iterative reinforcement improvement process, guided by feedback signals from evaluator models that compare generated documents to public and private reference corpora. The pipeline is structured around five main phases: (1) clinical entity keyword extraction using QuickUMLS \cite{soldaini2016quickumls} to ensure no direct patient identification; (2) supervised fine-tuning on pseudo-anonymized or differentially private documents; (3) candidate document generation; (4) evaluation by scorer models quantifying semantic quality; and (5) alignment through reinforcement learning using DPO \cite{rafailov2023dpo} or KTO \cite{ethayarajh2024kto} for continuous improvement.

\minitoc

\section{Extension of evaluation benchmarks}
Since our method shares several conceptual similarities with KnowledgeSG \cite{knowledgesg2024}, we sought to establish a comparison with this approach. KnowledgeSG proposes a client-server framework where the client possesses private knowledge (for example, a clinical model trained on sensitive data) while a knowledge distillation process occurs via a remote professional model. Like our pipeline, this method adopts an iterative approach and integrates alignment mechanisms and sample selection: it generates multiple candidates then progressively filters them via a reference professional model to select the best samples, thus countering the performance degradation induced by gradient noise when using differential privacy.

This methodological convergence naturally led us to adopt the AlpaCare benchmark \cite{zhang2023alpacare}, a doctor-patient conversational dataset in chatbot format, used both as a reference professional model in KnowledgeSG and as an evaluation benchmark in our context. This choice presents a dual advantage: on one hand, it enables direct comparison with KnowledgeSG on common evaluation ground, and on the other hand, it addresses a different type of task from our usual evaluations, thus broadening the validation spectrum of our approach. The evaluation follows a preference evaluation protocol where a judge compares pairwise the responses generated by our model with those from established reference models (text-davinci-003 \cite{OpenAI2022TextDavinci003}, GPT-3.5-turbo \cite{OpenAI2022GPT35Turbo}, GPT-4 \cite{OpenAI2024GPT4o}, and Claude-2 \cite{Anthropic2023Claude2}). To ensure result robustness, each sample is evaluated twice with swapped positions to avoid positional bias.

For our approach to adapt to this benchmark, we modified our method as follows: our SynthKG pipeline specifically generates synthetic patient questions from extracted medical keywords, while corresponding responses are produced by AlpaCare, which serves as the reference professional model. The process is structured as follows: we first generate clinical keywords, from which we synthesize patient questions, then AlpaCare generates the associated expert responses. This combination produces a question-answer pair dataset that is then used to fine-tune a model capable of answering questions from the AlpaCare dataset during evaluation.

\section{Quality-Oriented Synthetic Document Evaluation}

Given the results obtained with our initial pipeline, we identified potential improvement margins in the quality of synthetic document selection. While SemScore provided a foundation for semantic similarity assessment, we hypothesized that more sophisticated evaluation mechanisms could better capture the nuanced requirements of high-quality clinical synthetic documents. This section explores multi-faceted quality evaluation approaches that go beyond simple semantic similarity.

\subsection{Medical Domain Filtering}

The first quality enhancement we implemented was a medical domain filtering process adapted from KnowledgeSG \cite{knowledgesg2024}. This filter addresses a fundamental quality requirement: ensuring that generated documents remain medically relevant and do not drift into non-clinical domains during the generation process. The filtering mechanism operates by identifying whether a generated document or example constitutes a candidate whose subject matter falls within the medical domain.

This filtering step serves multiple purposes. First, it discards off-topic or non-medical generations that may arise from the stochastic nature of large language models, particularly when keyword guidance is sparse or ambiguous. Second, it ensures that synthetic documents maintain clinical relevance, a critical requirement for downstream medical NLP tasks. Third, it provides a coarse-grained quality gate before more expensive evaluation mechanisms are applied, reducing computational costs by filtering out obviously unsuitable candidates early in the pipeline.

The medical domain classifier was trained to distinguish clinical text from general domain text, focusing on lexical markers, medical terminology density, and structural patterns characteristic of clinical documentation. While simple in concept, this filter proved effective at eliminating generations that deviated significantly from the expected clinical context.

\subsection{LLM-as-Judge Evaluation}

Following medical domain filtering, we explored advanced LLM-as-judge evaluation techniques \cite{archit2024criticalevaluation,lee2024selfjudge}. This approach represents a paradigm shift from traditional automatic metrics by leveraging the contextual understanding and reasoning capabilities of large language models to evaluate document quality.

The LLM-as-judge framework operates by providing a large language model with a carefully designed prompt that instructs it to evaluate synthetic documents along multiple quality dimensions. Unlike SemScore, which measures semantic similarity through embedding distance, or traditional metrics like BLEU or ROUGE, which rely on surface-level token overlap, LLM-as-judge can assess higher-level qualities such as clinical coherence, factual plausibility, narrative flow, and adherence to medical documentation standards.

We designed evaluation prompts that asked judge models to score documents on a 1-5 scale based on clinical relevance, coherence, and appropriateness. The prompts explicitly instructed the model to consider whether the document could plausibly represent a real clinical note, whether the medical information presented was internally consistent, and whether the temporal sequence of events made clinical sense. For this evaluation, we selected a sample of one thousand documents that obtained a score of five according to KnowledgeSG's evaluation system as a reference point.

The advantage of LLM-as-judge lies in its ability to capture nuanced quality aspects that are difficult to formalize in traditional metrics. However, this approach introduces new challenges: judge models may exhibit biases, their evaluations may not perfectly align with human expert judgment, and the computational cost of LLM evaluation is substantially higher than traditional metrics. Furthermore, the reliability and calibration of LLM judges remains an active research question, particularly when evaluating specialized medical content.

\subsection{Composite Evaluation Strategy}

Our composite evaluation strategy combined medical filtering and LLM-as-judge in a two-stage pipeline. First, medical domain filtering eliminated clearly unsuitable candidates, reducing the evaluation pool. Second, LLM-as-judge provided fine-grained quality scoring on the filtered candidates. This hierarchical approach aimed to balance computational efficiency with evaluation sophistication.

We also explored weighted combination strategies where multiple quality signals—SemScore, medical domain probability, and LLM-judge scores—were aggregated to produce a final quality estimate. The intuition was that different metrics capture complementary aspects of quality: SemScore measures semantic fidelity to real documents, medical filtering ensures domain relevance, and LLM-as-judge assesses higher-level coherence and plausibility.

\subsection{Results and Analysis}

The detailed results presented in Table~\ref{tab:alpacare_detailed_results} reveal several important observations regarding our pipeline's performance. Analyzing first the results without medical filtering (using only SemScore), our SynthKG variants achieve respectable performance: the DPO-1 variant reaches 0.520 on average, while the SFT variant obtains 0.510. These scores position us competitively, significantly surpassing DP-Instruct-ICL (0.266), which according to the literature represents the best of traditional differential privacy methods.

The combined application of medical filter and LLM-as-judge (upper section of Table~\ref{tab:alpacare_detailed_results}) shows mixed results: DP-SFT + KTO reaches 0.539 and DP-SFT + DPO 0.525. While these more elaborate evaluation techniques yield slightly better results than the SemScore approach alone, the improvement remains limited and does not necessarily justify the additional complexity cost. The most notable gains are observed primarily against Text-davinci-003, which is the smallest of the evaluator models, while improvements remain marginal for other more sophisticated evaluation models. This performance does not surpass KnowledgeSG, which remains the reference method at 0.562, nor even AlpaCare which maintains its position at 0.536.

The limited magnitude of these improvements raises important questions about the cost-effectiveness of more complex evaluation approaches. We propose two hypotheses to explain these results. First, the bootstrap phase may be insufficient: our initial supervised fine-tuning on a small seed set may not provide enough diversity for the generator to benefit from more sophisticated selection mechanisms. Second, the evaluation techniques themselves may not provide sufficiently discriminative signals for this specific domain: medical conversational text may require domain-specific evaluation criteria that general-purpose LLM judges cannot adequately capture.

Our exchanges with the authors of KnowledgeSG also revealed that certain figures from the original paper should be taken with caution, identifying potential problems in how results were constructed. The comparison between filtered and unfiltered versions thus suggests that while quality-oriented evaluation improvements offer theoretical advantages, their practical benefits in our specific pipeline remain moderate, highlighting the need for further investigation into domain-specific evaluation mechanisms tailored to medical synthetic data generation.

\begin{table}[h]
\centering
\begin{tabular}{lccccc}
\toprule
\textbf{Method} & \textbf{Text-davinci-003} & \textbf{GPT-3.5-turbo} & \textbf{GPT-4} & \textbf{Claude-2} & \textbf{Average} \\
\midrule
\multicolumn{6}{l}{\textit{With medical filters and LLM-as-judge}} \\
\midrule
DP-SFT + DPO & 0.660 & 0.515 & 0.470 & 0.470 & \textbf{0.528} \\
DP-SFT + KTO & 0.659 & 0.512 & 0.492 & 0.490 & \textbf{0.540} \\
\midrule
\multicolumn{6}{l}{\textit{Without filters (SemScore only)}} \\
\midrule
DP-SFT + DPO & 0.592 & 0.504 & 0.478 & 0.468 & \textbf{0.513} \\
DP-SFT + KTO & 0.630 & 0.510 & 0.476 & 0.469 & \textbf{0.521} \\
\midrule
\multicolumn{6}{l}{\textit{Baselines and traditional DP methods}} \\
\midrule
AlpaCare & 0.666 & 0.506 & 0.474 & 0.497 & 0.536 \\
KnowledgeSG & \underline{\textbf{0.776}} & 0.530 & 0.457 & 0.488 & \underline{\textbf{0.562}} \\
DP-Instruct-ICL & 0.382 & 0.295 & 0.187 & 0.199 & 0.266 \\
Non-Private & 0.389 & 0.303 & 0.151 & 0.179 & 0.255 \\
\bottomrule
\end{tabular}
\caption{Detailed results on the AlpaCare benchmark by evaluator model (win rates). Our methods' averages are in bold. Best performances are underlined.}
\label{tab:alpacare_detailed_results}
\end{table}

\section{Privacy-Preserving Evaluation and Information Leakage Analysis}

While the previous section focused on improving document quality through more sophisticated evaluation mechanisms, a critical question remains: how can we ensure that the evaluation process itself does not compromise patient privacy? This section addresses the privacy implications of our evaluation pipeline and explores mechanisms to quantify and mitigate information leakage risks.

\subsection{Privacy Considerations in the Evaluation Pipeline}

Our evaluation architecture introduces a potential privacy vulnerability: the evaluator model must access real private documents to compare them with synthetic candidates and produce quality scores. Although only numerical scores cross the privacy boundary from the private to the public environment, we must rigorously analyze whether these scores inadvertently leak information about the private dataset.

The use of differential privacy during model training provides formal privacy guarantees, but comes at a cost. As noted in our experiments, differential privacy can introduce noise that degrades model performance, creating a fundamental tension between privacy preservation and synthetic document quality. Moreover, the use of differential privacy can unfortunately further noise the model, although it guarantees the absence of indistinguishable leaks. This privacy-utility trade-off becomes particularly acute in the medical domain, where both data protection and clinical accuracy are paramount.

\subsection{Threat Model and Attack Scenarios}

Given the critical privacy issues inherent to medical synthetic data generation, we developed a rigorous evaluation protocol to quantify re-identification risks of our generated documents. The central objective is to determine whether a synthetic document can be attributed to the specific patient whose data was used for its generation, constituting an unacceptable personal information leak in a medical context.

Our threat model fundamentally differs from traditional methodologies such as canary attacks used in KnowledgeSG. Canary attacks prove de facto inapplicable to our context since our method relies on keywords extracted from an ontological base (UMLS) that do not allow direct patient identification. These normalized keywords constitute by nature non-identifying elements, making the canary attack moot. Moreover, our threat model differs in a crucial way: the end user only accesses the synthetic dataset without being able to directly query the generative model, limiting the attack surface compared to interactive generation scenarios.

Instead, we focus on two complementary attack scenarios that better reflect realistic threats in our setting. First, we consider an adversary who possesses the synthetic dataset and attempts to infer which real patients contributed to the training data. Second, we examine whether an adversary with access to both synthetic and real documents can reliably link synthetic documents back to their source patients, even without explicit identifying information.

For the AlpaCare dataset, we had to circumvent a major methodological bias: conversational documents are not naturally attributed to unique patients. Prior clustering using BioLORD-2023-M allowed creating coherent pseudo-patients from 10,000 messages grouped into 4,000 user clusters, avoiding artifacts from initial grouping and making the evaluation more realistic and demanding.

\subsection{Re-identification Attack Implementation}

We implemented two complementary types of attacks to assess privacy vulnerabilities from different angles:

\paragraph{Linkage Attack (1-vs-N):} The linkage attack attempts to globally attribute each synthetic document to its source patient among the entire dataset. For each synthetic document, we compute similarity scores (using both TF-IDF and neural embeddings via all-mpnet-base-v2) against all real patient documents and predict the source as the highest-scoring candidate. The random baseline for this attack is near-zero (1/N where N is the number of patients), making any success rate above this threshold concerning from a privacy perspective.

\paragraph{Discrimination Attack (1-vs-2):} The discrimination attack performs a repeated binary discrimination test. For each synthetic document, we compare it three times to pairs composed of the real patient's private document versus that of a different randomly selected patient. The attacker must identify which of the two candidates is the true source. This attack has a theoretical baseline of 50\% (random guessing), and performance significantly above this threshold indicates that synthetic documents retain identifying characteristics.

The discrimination attack is particularly revealing because it represents a realistic scenario where an adversary has limited information—they know a synthetic document corresponds to one of two patients and must distinguish between them. Success at this task demonstrates that synthetic documents preserve patient-specific patterns even after the anonymization and synthesis process.

\subsection{Privacy Attack Results and Analysis}

The obtained results reveal contrasting vulnerabilities according to datasets and methods employed (see Table~\ref{tab:privacy_attacks}). For the linkage attack, AlpaCare presents re-identification rates of 49.5\% (TF-IDF) and 44.9\% (all-mpnet-base-v2), while MIMIC-III reaches more concerning levels of 76.0\% and 73.3\% respectively. The discrimination attack reveals even more critical vulnerabilities: AlpaCare reaches 96.6\% with all-mpnet-base-v2 and 89.1\% with TF-IDF, while MIMIC-III presents comparable results at 95.2\% and 94.1\%.

These results highlight a fundamental paradox of our approach. The linkage attack, which simulates an adversary attempting to identify a document's origin among all possible patients, presents moderate success rates (45-76\%). While these rates are substantially above the random baseline, they suggest that global re-identification is not trivial—an adversary cannot reliably pinpoint a patient's identity from the synthetic data alone without additional context.

In stark contrast, the discrimination attack achieves near-perfect performance (89-96\%). This disparity reveals a critical vulnerability: synthetic documents, while not trivially associable with their source in a global identification context, nevertheless retain characteristic fingerprints sufficiently distinctive to enable quasi-certain attribution in direct pairwise comparisons. The documents preserve patient-specific writing patterns, medical history signatures, or other latent features that allow reliable discrimination when the search space is constrained.

This observation raises important questions about the actual effectiveness of our privacy preservation method. The high discrimination attack success suggests that our keyword-guided generation approach, while preventing explicit PII leakage, may inadvertently preserve implicit patient identifiers through the selection and arrangement of medical concepts. The evaluator model, by scoring similarity to real documents, may be reinforcing these patient-specific patterns rather than eliminating them.

\subsection{Implications for Privacy-Preserving Evaluation}

These findings have direct implications for the design of privacy-preserving evaluation mechanisms. First, they demonstrate that even when only numerical scores cross the privacy boundary, the synthetic documents produced using these scores can leak information about the private dataset through their content and structure. Second, they highlight the need for explicit privacy-preserving mechanisms in the evaluation process itself, such as adding calibrated noise to scores or using differentially private training for the evaluator model.

The privacy-utility trade-off becomes clear: stronger privacy guarantees (e.g., higher differential privacy noise) would reduce re-identification risks but may also degrade the quality signal needed to improve synthetic document generation. Future work must explore this frontier more systematically, potentially through privacy budget allocation strategies that balance the privacy cost of evaluation against the utility gain from improved selection. Moreover, the necessity of developing more robust generation techniques to mitigate these vulnerabilities in terms of re-identification remains an open challenge, possibly through techniques such as adding controlled noise to keyword sequences or employing adversarial training to remove patient-specific signatures while preserving clinical validity.

\begin{table}[h]
\centering
\begin{tabular}{lcc}
\toprule
\textbf{Dataset} & \textbf{Attribution} & \textbf{Discrimination} \\
\midrule
AlpaCare (TF-IDF) & 0.495 & 0.891 \\
AlpaCare (Embeddings) & 0.449 & 0.966 \\
AlpaCare (Random) & - & 0.485 \\
MIMIC-III (TF-IDF) & 0.760 & 0.941 \\
MIMIC-III (Embeddings) & 0.733 & 0.952 \\
MIMIC-III (Random) & - & 0.499 \\
\bottomrule
\end{tabular}
\caption{Privacy attack results (DP-SFT+DPO with medical filters)}
\label{tab:privacy_attacks}
\end{table}
